\ulohahead{společná informatika \textnormal{\footnotesize[úroveň 2]}}{Dvouúrovňové stránkování a překlad adresy}
\begin{zadani}
16bitová virtuální adresa, velikost stránky 256\,B. Číslo stránky je rozděleno na dvě úrovně po 4 bitech. Virtuální adresa je \texttt{0x1234}.
\begin{enumerate}[label=(\alph*)]
  \item \bod{1} Rozdělte adresu na index L1, index L2 a offset.
  \item \bod{2} L1 tabulka v položce \texttt{0x1} ukazuje na L2 tabulku, jejíž položka \texttt{0x2} udává rámec \texttt{0x7}. Spočtěte fyzickou adresu.
  \item \bod{3} Proč se používá víceúrovňová tabulka místo jedné ploché? K čemu je TLB?
\end{enumerate}
\end{zadani}

\begin{reseni}
\textbf{(a)} Stránka 256\,B $\Rightarrow$ offset má $\log_2 256=8$ bitu. Zbývá 8 bitu čísla stránky, rozdělené na L1 (4 bity) a L2 (4 bity). \texttt{0x1234} $=$ \texttt{0001\,0010\,0011\,0100}$_2$: offset (dolních 8 bitu) $=$ \texttt{0x34}, číslo stránky (horních 8) $=$ \texttt{0x12}, tj.\ L1 $=$ \texttt{0x1}, L2 $=$ \texttt{0x2}.

\textbf{(b)} L1[\texttt{0x1}] $\to$ L2 tabulka; L2[\texttt{0x2}] $=$ rámec \texttt{0x7}. Fyzická adresa $=$ (rámec $\ll 8$) $\,|\,$ offset $=$ \texttt{0x7}\,$\ll8\ |\ $\texttt{0x34} $=$ \texttt{0x0734}.

\textbf{(c)} Jedna plochá tabulka by musela mít položku pro \emph{každou} stránku adresového prostoru (i nepoužité), což je u velkých prostoru obrovské. \emph{Víceúrovňová} tabulka alokuje L2 podtabulky jen pro skutečně použité oblasti (řídký adresový prostor) - šetří paměť. \emph{TLB} (translation lookaside buffer) je malá rychlá cache nedávných překladu stránka$\to$rámec, aby se nemusela při každém přístupu procházet tabulka v paměti.
\end{reseni}

\begin{jadro}
Adresa = (číslo stránky $\Vert$ offset); offset má $\log_2$(velikost stránky) bitu. Číslo stránky se u víceúrovňové tabulky dělí na indexy úrovní. Fyzická $=$ rámec $\ll$ (bity offsetu) $|$ offset. Víceúrovňová = šetří paměť u řídkého prostoru; TLB = cache překladu.
\end{jadro}

\kalibrace{Překlad adres je opakovaný výpočet (architektura/OS ~60\,\%); dvouúrovňová varianta je těžší než v úrovni 1. Rozdělení adresy + výpočet rámce = 2-3 body.}
\past{Offset má tolik bitu, kolik je $\log_2$(velikost stránky). Fyzická adresa skládá \emph{rámec} (ne číslo stránky) s offsetem. Bity čísla stránky se dělí mezi úrovně.}
